% TOP_PROBLEM: {"problemId": "problem.lovasz-conjecture-on-hamiltonian-paths-in-vertex-transitive-graphs", "canonicalTitle": "Lovász Conjecture on Hamiltonian Paths in Vertex-Transitive Graphs", "releaseRank": 436, "primaryDomain": "combinatorics_discrete_geometry", "primaryDomainLabel": "Combinatorics & discrete geometry", "displayStatus": "open_disputed_claim", "releaseStatus": "open_disputed_claim"}
% !TeX program = lualatex
% ATTEMPT_STATUS: unresolved
\documentclass[11pt]{article}
\usepackage[margin=25mm]{geometry}
\usepackage{fontspec}
\setmainfont{Latin Modern Roman}
\usepackage{amsmath,amssymb,mathtools}
\usepackage{unicode-math}
\setmathfont{Latin Modern Math}
\usepackage{xurl,hyperref}
\hypersetup{colorlinks=true,urlcolor=blue}
\setcounter{secnumdepth}{0}
\setlength{\parindent}{0pt}
\setlength{\parskip}{5pt}
\setlength{\emergencystretch}{3em}
\begin{document}
\section*{\texorpdfstring{Lovász Conjecture on Hamiltonian Paths in Vertex-Transitive Graphs}{Lovász Conjecture on Hamiltonian Paths in Vertex-Transitive Graphs}}\label{436}
\subsection{Definitions and mathematical statement}
Let $G=(V,E)$ be a finite connected simple graph. It is vertex-transitive if for every $u,v\in V$ an adjacency-preserving bijection of $V$ maps $u$ to $v$. A Hamiltonian path is an ordering $v_1,\ldots,v_{|V|}$ of all vertices with $v_iv_{i+1}\in E$ for every adjacent pair in the ordering. The conjecture is
\[
 G\text{ connected and vertex-transitive}
 \quad\Longrightarrow\quad G\text{ has a Hamiltonian path}.
\]
The endpoints need not be adjacent. Thus the target is a path, not the stronger assertion of a Hamiltonian cycle. The one-vertex graph is allowed with its trivial path.
\par\smallskip\textit{Formulation and concept references:} \href{https://arxiv.org/abs/2606.09742}{[S1]}.

\subsection{Short English statement}
Must every finite connected graph that looks the same at every vertex have a path visiting every vertex exactly once?
\subsection{Research attempt}
\paragraph{Scope, process, and first gap.}
This attempt by GPT-6 Astra Ultra was prepared on 27 September 2026. It proves the conjecture for connected Cayley graphs of finite abelian groups by an explicit coset traversal, checks elementary consequences of vertex transitivity, and uses the Petersen graph to test a proposed upgrade from paths to cycles. These are established partial arguments, not a new proof of Lovász's conjecture.

The catalogue's disputed-claim status is retained. The primary source [S1] concerns progress towards the general conjecture; this attempt does not adopt a claimed complete proof. The first missing step is extending the compatible coset traversal to arbitrary vertex-transitive graphs. Such graphs need not come with a regular abelian group action, and an arbitrary transitive automorphism action does not supply the needed labels or endpoint alignment.

\paragraph{Basic consequences that fall short of Hamiltonicity.}
A vertex-transitive graph is regular, because an automorphism maps the neighbor set of one vertex bijectively to that of another. A connected regular graph of degree zero has one vertex, and degree one gives $K_2$. Both have the required path. Degree two gives a cycle, also with a Hamiltonian path after deleting one edge.

For a connected vertex-transitive graph with at least two vertices, no vertex is a cut vertex. Indeed, take a spanning tree and one of its leaves. Removing that leaf leaves the rest of the spanning tree connected, so it is not a cut vertex of the original graph. Being a cut vertex is invariant under automorphisms; transitivity then excludes every cut vertex.

These facts eliminate simple obstructions such as a branching articulation point. They do not create a spanning path. A proof that merely upgrades connectivity to absence of cut vertices would leave most of the problem untouched.

\paragraph{An explicit abelian Cayley construction.}
Let $A$ be a finite abelian group, written additively, and let $S=-S$ be a generating set not containing zero. Its undirected Cayley graph has edges $a$--$(a+s)$ for $s\in S$. Translations are automorphisms, so it is vertex-transitive, and generation is exactly connectedness.

Choose generators $s_1,\ldots,s_t$ from $S$ successively so that
\[
 H_0=\{0\}\subsetneq H_1\subsetneq\cdots\subsetneq H_t=A,
 \qquad H_i=\langle s_1,\ldots,s_i\rangle.
\]
We construct a Hamiltonian path in each $H_i$ using only these generators. Begin with the trivial path in $H_0$. Suppose
$P=(p_1,\ldots,p_h)$ visits $H_{i-1}$ once. The quotient
$H_i/H_{i-1}$ is cyclic, generated by $s_i+H_{i-1}$; write its order as $r$.

Its distinct cosets are $H_{i-1}+js_i$, $0\le j<r$. Traverse them in that order, using
\[
 (p_1+js_i,\ldots,p_h+js_i)\quad\text{if }j\text{ is even},
\]
and the reversed order if $j$ is odd. Each row is a path because translation preserves differences in $S$, and reversing a path is allowed because the graph is undirected. Consecutive rows meet at corresponding endpoints: the last element of row $j$ and first element of row $j+1$ differ by $s_i$. Hence an allowed edge joins them.

Distinct rows are disjoint cosets, and each row lists all of its coset. The concatenation therefore visits every element of $H_i$ once. Induction yields a Hamiltonian path in the original graph. This is a complete proof for every finite abelian Cayley graph, without an assumption that one generator alone has full order.

\paragraph{What the construction actually uses.}
The alternating reversal is essential. If every row used the same orientation, the joining difference would be $p_1-p_h+s_i$, which need not lie in $S$. In the construction it is exactly $s_i$, because consecutive rows expose the same endpoint of $P$.

Likewise the quotient is a genuine group and its cosets can be translated while preserving the generating edges. For nonabelian groups, left and right translation conventions must be tracked and subgroup normality can fail. For a general vertex-transitive graph, identifying vertices with arbitrary automorphisms repeats vertices whenever stabilizers are nontrivial. None of these obstacles is removed by the existence of a transitive action alone.

This clarifies the first gap in a block-induction strategy: one must produce spanning paths inside blocks with endpoints compatible with the available inter-block edges. A Hamiltonian path in a quotient graph plus unrelated Hamiltonian paths in its blocks does not guarantee that compatibility.

\paragraph{The Petersen graph: a path need not close.}
Use outer vertices $v_0,\ldots,v_4$ and inner vertices $u_0,\ldots,u_4$, with indices modulo five. Edges are
\[
 v_iv_{i+1},\qquad v_iu_i,\qquad u_iu_{i+2}.
\]
This is the Petersen graph. An explicit Hamiltonian path is
\[
 v_0,v_1,v_2,v_3,v_4,u_4,u_1,u_3,u_0,u_2.
\]
Every listed consecutive pair is one of the specified edges, and all ten vertices occur once.

For completeness, it is vertex-transitive via its equivalent description as the graph on two-element subsets of a five-element set, adjacent exactly when disjoint. Every permutation of the ground set preserves disjointness, and any two two-element subsets are carried to one another. This graph is isomorphic to the displayed outer/inner presentation.

It has no Hamiltonian cycle. Such a cycle would cross between the outer and inner five-vertex sets an even positive number of times, so it would use two or four spokes. If it used two, the selected outer edges would form a spanning path of the outer 5-cycle; its endpoints have indices differing by one. The selected inner edges would form a spanning path of the inner 5-cycle, whose endpoints have indices differing by two. The two spoke indices cannot satisfy both requirements.

If it used four spokes, rotate labels so that the omitted spoke is $v_0u_0$. At $v_0$ both outer edges must be selected, while every other outer vertex has one selected outer edge. This forces outer edges
$v_4v_0,v_0v_1,v_2v_3$. The same degree count inside forces
$u_0u_2,u_0u_3,u_1u_4$. Together with the four spokes these form two separate 5-cycles, not one spanning cycle. Thus four spokes also fail.

The example invalidates any proof step that assumes a longest spanning path can always have its endpoints joined. It supports the distinction in the statement; it does not contradict its Hamiltonian-path conclusion.

\paragraph{Verification and outcome.}
The finite checks implement the alternating coset traversal for several products of cyclic groups and verify uniqueness and adjacency of every listed vertex. Separate exact searches verify the displayed Petersen path and the absence of a Hamiltonian cycle. The written cycle proof explains the result without relying on enumeration.

\textbf{UNRESOLVED.} The conjecture is proved here for finite abelian Cayley graphs and the elementary low-degree cases. The endpoint compatibility required for a general block or automorphism argument remains unproved. No counterexample to the path conjecture is produced, and the stronger cycle assertion is explicitly rejected.

\subsection{Sources}
\begin{itemize}
\item[S1]Bucic, Christoph, Pokrovskiy, and Steiner, Towards the Lovasz conjecture via sublinear expanders (2026), Abstract.
\url{https://arxiv.org/abs/2606.09742}.
\textit{Formulation source.}
\end{itemize}
\end{document}
